%%
%% Copyright 2007, 2008, 2009 Elsevier Ltd
%%
%% This file is part of the 'Elsarticle Bundle'.
%% ---------------------------------------------
%%
%% It may be distributed under the conditions of the LaTeX Project Public
%% License, either version 1.2 of this license or (at your option) any
%% later version.  The latest version of this license is in
%%    http://www.latex-project.org/lppl.txt
%% and version 1.2 or later is part of all distributions of LaTeX
%% version 1999/12/01 or later.
%%
%% The list of all files belonging to the 'Elsarticle Bundle' is
%% given in the file `manifest.txt'.
%%

%% Template article for Elsevier's document class `elsarticle'
%% with numbered style bibliographic references
%% SP 2008/03/01
%%
%%
%%
%% $Id: elsarticle-template-num.tex 4 2009-10-24 08:22:58Z rishi $
%%
%%
%\documentclass[preprint,12pt,3p]{elsarticle}

%% Use the option review to obtain double line spacing
%% \documentclass[preprint,review,12pt]{elsarticle}

%% Use the options 1p,twocolumn; 3p; 3p,twocolumn; 5p; or 5p,twocolumn
%% for a journal layout:
%% \documentclass[final,1p,times]{elsarticle}
%% \documentclass[final,1p,times,twocolumn]{elsarticle}
%% \documentclass[final,3p,times]{elsarticle}
%% \documentclass[final,3p,times,twocolumn]{elsarticle}
%% \documentclass[final,5p,times]{elsarticle}
\documentclass[final,5p,times,twocolumn]{elsarticle}

%% if you use PostScript figures in your article
%% use the graphics package for simple commands
%% \usepackage{graphics}
%% or use the graphicx package for more complicated commands
%% \usepackage{graphicx}
%% or use the epsfig package if you prefer to use the old commands
%% \usepackage{epsfig}

%% The amssymb package provides various useful mathematical symbols
\usepackage{amssymb}
%% The amsthm package provides extended theorem environments
%% \usepackage{amsthm}

%% The lineno packages adds line numbers. Start line numbering with
%% \begin{linenumbers}, end it with \end{linenumbers}. Or switch it on
%% for the whole article with \linenumbers after \end{frontmatter}.
\usepackage{lineno}

%% natbib.sty is loaded by default. However, natbib options can be
%% provided with \biboptions{...} command. Following options are
%% valid:

%%   round  -  round parentheses are used (default)
%%   square -  square brackets are used   [option]
%%   curly  -  curly braces are used      {option}
%%   angle  -  angle brackets are used    <option>
%%   semicolon  -  multiple citations separated by semi-colon
%%   colon  - same as semicolon, an earlier confusion
%%   comma  -  separated by comma
%%   numbers-  selects numerical citations
%%   super  -  numerical citations as superscripts
%%   sort   -  sorts multiple citations according to order in ref. list
%%   sort&compress   -  like sort, but also compresses numerical citations
%%   compress - compresses without sorting
%%
%% \biboptions{comma,round}

% \biboptions{}

\usepackage{graphicx}
\usepackage[space]{grffile}
\usepackage{latexsym}
\usepackage{textcomp}
\usepackage{longtable}
\usepackage{tabulary}
\usepackage{booktabs,array,multirow}
\usepackage{amsfonts,amsmath,amssymb}
\providecommand\citet{\cite}
\providecommand\citep{\cite}
\providecommand\citealt{\cite}
\usepackage{url}
\usepackage{hyperref}
\hypersetup{colorlinks=false,pdfborder={0 0 0}}
\usepackage{etoolbox}
\makeatletter
\patchcmd\@combinedblfloats{\box\@outputbox}{\unvbox\@outputbox}{}{%
  \errmessage{\noexpand\@combinedblfloats could not be patched}%
}%
\makeatother
% You can conditionalize code for latexml or normal latex using this.
\newif\iflatexml\latexmlfalse
\providecommand{\tightlist}{\setlength{\itemsep}{0pt}\setlength{\parskip}{0pt}}%

\AtBeginDocument{\DeclareGraphicsExtensions{.pdf,.PDF,.eps,.EPS,.png,.PNG,.tif,.TIF,.jpg,.JPG,.jpeg,.JPEG}}

\usepackage[utf8]{inputenc}
\usepackage[english]{babel}



% Authorea add-on: we want to allow long links in a two-column style, we'll need to patch some floats.
\usepackage{etoolbox}
\makeatletter
\patchcmd\@combinedblfloats{\box\@outputbox}{\unvbox\@outputbox}{}{%
  \errmessage{\noexpand\@combinedblfloats could not be patched}%
}%
% Authorea add-on: ensure doi defined
\makeatletter
\providecommand{\doi}[1]{%
  \begingroup
    \let\bibinfo\@secondoftwo
    \urlstyle{rm}%
    \href{http://dx.doi.org/#1}{%
      doi:\discretionary{}{}{}%
      \nolinkurl{#1}%
    }%
  \endgroup
}

\def\ps@pprintTitle{%
 \let\@oddhead\@empty
 \let\@evenhead\@empty
 \def\@oddfoot{\selectlanguage{english}\footnotesize\itshape
       Preprint submitted to \ifx\@journal\@empty Elsevier
       \else\@journal\fi\hfill\today}%
 \let\@evenfoot\@oddfoot}

\makeatother

\journal{Elsevier Journal}

\begin{document}
\begin{frontmatter}

%% Title, authors and addresses

%% use the tnoteref command within \title for footnotes;
%% use the tnotetext command for the associated footnote;
%% use the fnref command within \author or \address for footnotes;
%% use the fntext command for the associated footnote;
%% use the corref command within \author for corresponding author footnotes;
%% use the cortext command for the associated footnote;
%% use the ead command for the email address,
%% and the form \ead[url] for the home page:
%%
%% \title{Title\tnoteref{label1}}
%% \tnotetext[label1]{}
%% \author{Name\corref{cor1}\fnref{label2}}
%% \ead{email address}
%% \ead[url]{home page}
%% \fntext[label2]{}
%% \cortext[cor1]{}
%% \address{Address\fnref{label3}}
%% \fntext[label3]{}

%% Patch date to be in main language
\title{Titan Haze Paper}%


\author[add1]{cihat kurt}%
\address[add1]{Texas A\&M University}%





%\begin{keyword}
%% keywords here, in the form: keyword \sep keyword
%% MSC codes here, in the form: \MSC code \sep code
%% or \MSC[2008] code \sep code (2000 is the default)
%\end{keyword}
\end{frontmatter}
\linenumbers

\section*{The Final Report}

{\label{the-final-report}}

I searched for the best-fitting phase functions among
nearly\(5\times 10^{5}\)model runs using an automated algorithm. The
algorithm determines different constraints for each altitude and finds
two parameter ranges for each wavelength channel (one for each
altitude).~It obtains the constraints from the maximum polarization, the
phase function at the forward scattering angles and the single
scattering albedo (SSA or \(\omega_{0}\)) by comparing the model
aggregates to the observations. Then based on the mean square error
(MSE) the algorithm decides whether a model aggregate fits to the
observations. As for the values of the phase functions at backscattering
angles (\(\nu\geq 75\ \)), the code checks the MSE and makes two
sample Anderson-Darling (AD) test to compare the curves of phase
functions at these angles. Then based on the MSE and the p-values of the
AD test the model aggregates are either decided to match the DISR phase
function or not. AD test provides a quantitative \textbf{measure} for
how much the shapes of the phase functions being compared resemble each
other as well as how close the individual points on the curves in
average.

\textbf{Even though I adopted a single phase function throughout the
atmosphere for each channel, SSA observations provide constraints on the
refractive index at two altitudes} . Therefore, for each layer, a unique
parameter ranges can be found by simply selecting the model aggregates
which agree with all three observations (SSA, phase function, linear
polarization). Besides, the phase function also constraints the
refractive index.\selectlanguage{english}
\begin{figure*}[h!]
\begin{center}
\includegraphics[width=0.98\columnwidth]{figures/image1/image1}
\caption{{Couldn't find a caption, edit here to supply one.
{\label{div-331178}}%
}}
\end{center}
\end{figure*}

As shown in \textbf{Figure x} , \textbf{Figure} \textbf{y and Figure
z}while \(n_{r}\) is mostly constrained by the phase function,
\(n_{i}\)is constrained mostly by SSA except that the phase
function in the blue channel also strongly correlated with
\(n_{i}\), suggesting strong absorption.\selectlanguage{english}
\begin{figure*}[h!]
\begin{center}
\includegraphics[width=0.98\columnwidth]{figures/image2/image2}
\caption{{Couldn't find a caption, edit here to supply one.
{\label{div-734570}}%
}}
\end{center}
\end{figure*}

The constraints on the refractive index from the phase functions agree
with the constraints from SSA except that blue-phase function suggests
larger \(n_{i}\) than the SSA measurement does in the blue
channel 80 km altitude. The algorithm finds almost no model that passes
all the filters at 80 km in the Blue channel without allowing at least
10\% error in SSA or more than 20\% error in the phase function at the
backscattering angles. This could be due to the error in SSA observation
or the blue phase function at this level given in the DISR model. In the
DISR model, both in the original (Tomasko et al., 2008)~and revised
version (Doose et al., 2016), blue phase function is constant throughout
the atmosphere while the SSA is changing. It has been suggested
by~Tomasko et al. (2008)~and Doose et al. (2016) that some condensation
at the aerosol surface occurs in the lower atmosphere resulting in a
decrease in absorption while keeping the phase function relatively
constant.

On the other hand, the fact that our model does not produce any aerosol
matching all the constraints within 5\% accuracy might suggest that
constant phase functions throughout the atmosphere might not be an
accurate representation of reality. Likewise, the SSA measurements at 80
km strongly suggest that the \(n_{i}\) values are considerably
smaller (\textbf{Figure xx)} than what the phase function suggests in
the blue channel (\textbf{Figure 1} ). Perhaps, the condensing particles
on the aggregate's surface not only change the \(w_{0}\) but
also the phase function of the haze particles. (\textbf{Give an example
from brown and black carbon, freshly produced and aged one).} However,
It is still possible that the new parameterization might also be the
source of this apparent disagreement, \textbf{or the SSA measurement may
not be accurate at 80 km in the blue channel.}\selectlanguage{english}
\begin{figure*}[h!]
\begin{center}
\includegraphics[width=0.98\columnwidth]{figures/image3/image3}
\caption{{Couldn't find a caption, edit here to supply one.
{\label{div-485060}}%
}}
\end{center}
\end{figure*}

Due to this \textbf{disagreement between the constraints from the phase
function and the SSA measurements in the blue channel, at 80 Km
altitude,} I applied a slightly different method to determine the
constraints for this particular case. In this method, the algorithm
first finds the monomer number and the size ranges from the DISR phase
functions and the maximum polarization without yet comparing
the\(w_{0}\) values. Next, inside the aggregate size range
(\(\text{N\ and\ }R_{m}\)) it determines the range of refractive index from
the \(\omega_{0}\ \)measurements within the 5\% error limit. As a
result, the constraints from \(\omega_{0}\) is semi-independent from
the backscattering values of the phase function but not from the values
of the forward angles of the phase function and the maximum
polarization, which are still set to fit within 5\% limit as they are
the only direct measurements apart from \(\omega_{0}\).

Table 1 and figure 2 \& 3 summarizes the parameters of the haze
particles. Constraints from the same channel provide the same parameter
range at each altitude except for the imaginary part of the refractive
index. This is consistent with the observation as the SSA measurements
imply that more absorbing particles are found at higher altitudes while
the size of the aggregates is somewhat constant.

Monomer size ranges found from each channel are consistent and about the
same at each altitude. On the other hand, at a given altitude, the
ranges of monomer number (\(N\)) differ depending on which
channel they are calculated from. As seen in \textbf{Figure x and Figure
y} , the range of monomer numbers in aggregates that fits the
observations is shifted upward in the Red channel comparing to those
obtained in the blue channel. The similar inconsistency between each
channel also appears in the overall size of the haze model particles
(\(R_{\text{agg}}),\) especially for the 80 km altitude.\selectlanguage{english}
\begin{longtable}[]{@{}lllllll@{}}
\toprule
\textbf{Altitude} & \textbf{Channel} & \(\mathbf{N}\) &
\(\mathbf{R}_{\mathbf{m}}\) & \(\mathbf{R}_{\mathbf{\text{agg}}}\) & \(\mathbf{n}_{\mathbf{r}}^{\mathbf{*}}\) &
\(\mathbf{n}_{\mathbf{i}}^{\mathbf{*}}\)\tabularnewline
\midrule
\endhead
\textbf{200 km} & \textbf{Red} & \(2950-4700\) & \(0.036-0.045\)
& \(0.61-0.65\) & \(1.62-1.83\) &
\(0.004-0.022\)\tabularnewline
& \textbf{Blue} & \(2200-3450\) & \(0.038-0.049\) &
\(0.56-0.65\) & \(1.77-2.0\) &
\(0.096-0.126\)\tabularnewline
\textbf{80} & \textbf{Red} & \(2950-4700\) & \(0.036-0.045\) &
\(0.61-0.65\) & \(1.62-1.83\) &
\(0.0005-0.008\)\tabularnewline
& \textbf{Blue} & \(2200-3450\) & \(0.036-0.048\) &
\(0.55-0.62\) & \(1.79-2.0\) &
\(0.01-0.066\)\tabularnewline
\bottomrule
\end{longtable}

Table 1 shows the results from all the constraints at each layer and
each channel. These results obtained by finding the model-aerosols
within 5\% accuracy to all the observations (Phase function, maximum
polarization, SSA) except at the backscattering an

Ghfh\selectlanguage{english}
\begin{figure*}[h!]
\begin{center}
\includegraphics[width=0.98\columnwidth]{figures/image4/image4}
\caption{{Couldn't find a caption, edit here to supply one.
{\label{div-576244}}%
}}
\end{center}
\end{figure*}

Figure 4\selectlanguage{english}
\begin{figure*}[h!]
\begin{center}
\includegraphics[width=0.98\columnwidth]{figures/image5/image5}
\caption{{Couldn't find a caption, edit here to supply one.
{\label{div-262578}}%
}}
\end{center}
\end{figure*}

Figure 5

\textbf{Table 2} shows the parameters of the haze particle obtained by
combining the constraints from both channels except for the refractive
index. Two different parameter ranges are calculated for each altitude.
One from combining only the intersection of the parameter ranges of the
blue and red channel and another is the union of two sets from each
channel, labeled as \emph{Inner} and \emph{Outer,} respectively,
in\textbf{Table 2} .

Therefore, if we neglect the model aggregates which do not agree with
the constraints from both channels (\emph{Inner combination),} then we
find that the haze particles have about \(3200\pm 250\ \)monomers each
have a radius of \(R_{m}=0.042\pm 0.004\). Otherwise accounting all the model
aggregates which agree with the constraints from at least one channel
(outer combination) suggest that the haze particles have
about\(N=3400\pm 1300\) monomers with monomer radius\(R_{m}=0.042\pm 0.007\).\selectlanguage{english}
\begin{longtable}[]{@{}lllll@{}}
\toprule
Combination & Altitude & N & \(\mathbf{R}_{\mathbf{m}}\) &
\(\mathbf{R}_{\mathbf{\text{agg}}}\)\tabularnewline
\midrule
\endhead
Inner & 200 km & \(2950\leq N\leq\ 3450\) & \(0.038\leq R_{m}\leq\ 0.045\) &
\(0.61\leq R_{\text{agg}}\leq 0.65\)\tabularnewline
& 80 km & \(2950\leq N\leq\ 3450\) & \(0.036\leq R_{m}\leq\ 0.045\) &
\(0.61\leq R_{\text{agg}}\leq 0.62\)\tabularnewline
Outer & 200 km & \(2200\leq N\leq\ 4700\) & \(0.036\leq R_{m}\leq\ 0.049\) &
\(0.56\leq R_{\text{agg}}\leq 0.65\)\tabularnewline
& 80 km & \(2200\leq N\leq\ 4700\) & \(0.036\leq R_{m}\leq\ 0.048\) &
\(0.55\leq R_{\text{agg}}\leq 0.65\)\tabularnewline
\bottomrule
\end{longtable}

Table 2 Parameter of Haze related to its size from the combination of
constraints in the SA channels. Inner means the parameter ranges
combined by the intersection of each channel's parameter ranges. The
outer refers to the combination by the outer edges.

Figure 5 shows the collection of all the fitting phase functions of the
model aggregates within the parameter range given in table 1. All the
phase functions that make up the table 2 aggregates have an excellent
agreement with the DISR model at all scattering angles for the Red phase
function except near the backscattering peak. As for the blue phase
functions, there is an apparent bias in the intermediate angles. Perhaps
the new parameterization does not extrapolate well at these angles and,
therefore, cause such a bias. On the other hand, considering that the
DISR phase functions were obtained from the old parameterization model
(Tomasko et al., 2008), It is also possible that DISR blue phase
function may not be accurate at the intermediate angles, albeit they
were also constrained by radiative transfer model to fit the
observations. Regardless, the algorithm did not use these intermediate
angles in the search for best-fitting parameters.\selectlanguage{english}
\begin{figure*}[h!]
\begin{center}
\includegraphics[width=0.98\columnwidth]{figures/image6/image6}
\caption{{Couldn't find a caption, edit here to supply one.
{\label{div-800293}}%
}}
\end{center}
\end{figure*}

Among the model runs shown in table 1 and \textbf{fig 3} the best
aggregate which fits the DISR phase functions, maximum polarization and
SSA measurements (above 80km) within \%5 error limit
has\(\sim 3.3\times 10^{3}\) monomers with \(R_{m}\sim\ 0.04\). As shown in
\textbf{fig below,} such an aggregate fits the DISR phase functions very
well in both channels.\selectlanguage{english}
\begin{figure*}[h!]
\begin{center}
\includegraphics[width=0.98\columnwidth]{figures/image7/image7}
\caption{{Couldn't find a caption, edit here to supply one.
{\label{div-630384}}%
}}
\end{center}
\end{figure*}

Similarly, When the model aggregates are compared to DISR phase
functions at each wavelength independently, the best-fitting phase
function at each wavelength agrees better with the DISR phase function
at the corresponding channel. However, there is a somewhat 20\%
difference in the size of the aggregates constrained by each
channel.\textbf{Figure 9} shows the phase functions of these model
aggregates in comparison to the DISR model.\selectlanguage{english}
\begin{figure*}[h!]
\begin{center}
\includegraphics[width=0.98\columnwidth]{figures/image8/image8}
\caption{{Couldn't find a caption, edit here to supply one.
{\label{div-637372}}%
}}
\end{center}
\end{figure*}

\textbf{Table 3} compares the parameters

nge between 1500 -- 6000.\selectlanguage{english}
\begin{longtable}[]{@{}llll@{}}
\toprule
\begin{minipage}[b]{0.22\columnwidth}\raggedright\strut
\strut
\end{minipage} & \begin{minipage}[b]{0.22\columnwidth}\raggedright\strut
N\strut
\end{minipage} & \begin{minipage}[b]{0.22\columnwidth}\raggedright\strut
\(\mathbf{R}_{\mathbf{m}}\ (\mu m)\)\strut
\end{minipage} & \begin{minipage}[b]{0.22\columnwidth}\raggedright\strut
\(\mathbf{R}_{\mathbf{\text{agg}}}\ (\mu m)\)\strut
\end{minipage}\tabularnewline
\midrule
\endhead
\begin{minipage}[t]{0.24\columnwidth}\raggedright\strut
This Work\strut
\end{minipage} & \begin{minipage}[t]{0.24\columnwidth}\raggedright\strut
\[\ \mathbf{\ 3300}\ \] \[2200\leq N\leq\ 4700\ \]\strut
\end{minipage} & \begin{minipage}[t]{0.24\columnwidth}\raggedright\strut
\[\mathbf{0.04}\mathbf{1}\] \[\ 0.038\leq R_{m}\leq\ 0.045\]\strut
\end{minipage} & \begin{minipage}[t]{0.24\columnwidth}\raggedright\strut
\[\mathbf{0.6}\mathbf{3}\] \[\ 0.61\leq R_{\text{agg}}\leq 0.65\]\strut
\end{minipage}\tabularnewline
\begin{minipage}[t]{0.24\columnwidth}\raggedright\strut
Tomasko et al. (2008)\strut
\end{minipage} & \begin{minipage}[t]{0.24\columnwidth}\raggedright\strut
\[\ \mathbf{3000\ }\] \[1500\leq N\leq\ 6000\]\strut
\end{minipage} & \begin{minipage}[t]{0.24\columnwidth}\raggedright\strut
\[\mathbf{0.05\ }\] \[0.03\leq R_{m}\leq 0.07\]\strut
\end{minipage} & \begin{minipage}[t]{0.24\columnwidth}\raggedright\strut
\[\ \mathbf{0.72}\]\strut
\end{minipage}\tabularnewline
\begin{minipage}[t]{0.24\columnwidth}\raggedright\strut
Tomasko et al. (2009)\strut
\end{minipage} & \begin{minipage}[t]{0.24\columnwidth}\raggedright\strut
\[\ \mathbf{4300}\ \] \[3000\leq N\leq 7000\]\strut
\end{minipage} & \begin{minipage}[t]{0.24\columnwidth}\raggedright\strut
\[\mathbf{0.04}\] \[0.03\leq R_{m}\leq 0.05\]\strut
\end{minipage} & \begin{minipage}[t]{0.24\columnwidth}\raggedright\strut
\[\ \mathbf{0.65}\] \(0.57\leq R_{\text{agg}}\leq 0.72\) *\strut
\end{minipage}\tabularnewline
\begin{minipage}[t]{0.22\columnwidth}\raggedright\strut
Other\strut
\end{minipage} & \begin{minipage}[t]{0.22\columnwidth}\raggedright\strut
\strut
\end{minipage} & \begin{minipage}[t]{0.22\columnwidth}\raggedright\strut
\strut
\end{minipage} & \begin{minipage}[t]{0.22\columnwidth}\raggedright\strut
\strut
\end{minipage}\tabularnewline
\bottomrule
\end{longtable}

Table 3 Comparison of the Titan's haze parameters estimated by different
studies. The best estimate given in each work is shown as bold text. The
range of the parameter estimates is given if available. * Aggregate size
is not given explicitly in the reference text but calculated based on
given monomer size uncertainty and the corresponding monomer numbers
estimate.

( From Tomasko 2009: ``Different authors have found the monomer size of
Titan's haze aerosols to be due to different effects. Cabane et al.
(1992) and Ran- nou et al. (2000) found that the monomer size depended
most on the altitude of formation of the aggregate particles. In their
models, larger monomers could be supported at lower altitudes, and this
determined the monomer size of the aggregate aerosols. In a new study,
Rodin et al. (2009) used a different way of estimating the electrical
charge on the aerosol particles. They find that monomers smaller than
about 0.05 lm'')

Because of the possibility of some condensation on the surface of the
haze particles below 80 km, which in turn reduces \(n_{i}\),
here I only compare the refractive index constrained by the measurements
at 200-km altitude to the laboratory studies. The ranges of refractive
index at 491 nm and 935 nm wavelength obtained in this work are closer
to the ones obtained by Khare et al. (1984) than the other laboratory
studies shown in \textbf{Figure 7} . At 491 nm \(n_{r}\) is
found to be in between 1.77 -- 2.0 and \(n_{i}\) is in between
0.-096 -- 0.126 while Khare et al. (1984) found that \(n_{r}\)=
\textbf{XXX} and \(n_{i}\)=\textbf{YYY} . At 935 nm in this
work \(n_{r}\ \)is between 1.62 -- 1.83 and\(\ n_{i}\)
ranges from 0.004 to 0.022 whereas Khare et al. (1984) found that
\(n_{r}=xx\) and \(n_{i}=yy\).\selectlanguage{english}
\begin{figure*}[h!]
\begin{center}
\includegraphics[width=0.98\columnwidth]{figures/image9/image9}
\caption{{Couldn't find a caption, edit here to supply one.
{\label{div-628811}}%
}}
\end{center}
\end{figure*}

\subsection*{Conclusion}

{\label{conclusion}}

DISR model for Titan's haze was published in (Tomasko et al., 2008). The
model includes phase functions for several wavelengths, but the ones for
491 nm and 935 nm are the focus of this work because these are the peak
wavelengths that were detected by the Solar Aureole camera, which
directly measured the forward scattering of solar radiation. Doose et
al. (2016) revised the phase functions and the single scattering albedo
measurements based on more detailed analyses of measurements inside and
outside of the Titan's atmosphere. As opposed to the original DISR model
(Tomasko et al., 2008), they used a single phase function throughout the
atmosphere for each wavelength. I adopted the DISR model phase functions
and SSA from Doose et al. (2016). As for the linear polarization,
Tomasko et al., (2009) analyzed the polarization data in more detail
than Tomasko et al., 2008 did. I, therefore, utilize the polarization
data given by Tomasko et al., (2009) in this work.

Since the DISR measurements are analyzed in blue and red channels, for
each altitude, two sets of constraints on the parameters of haze are
obtained. While the constraints on the refractive index are unique to
each wavelength, the monomer number and monomer size are supposed to be
the same no matter which channel's data is used to constraint them.
Unfortunately, there is some discrepancy between the ranges of N
obtained from each channel (see \textbf{table 1 figure 2 and figure 3}
). For this reason, the combination of the constraints from both
channels results in more extensive parameter ranges than the ones
obtained by each channel, unless only the intersections of the parameter
ranges are considered to match the measurements (see \textbf{Table 2} ).
\textbf{In any case} , the parameter ranges are in line with, but
narrower than, the ones obtained by earlier works (\textbf{see Table 3)}

Based on all the constraints from DISR measurements, the best fitting
aggregates contain 2900-3500 monomers with monomer size,
\(R_{m}\) is between 0.038 -- 0.045. The aggregates with larger
N have usually smaller \(R_{m}\) or vice versa. Hence, the
range of the overall size of the aggregates which fits the observations
is narrower:\(0.61\leq R_{\text{agg}}\leq 0.63\), where \(R_{\text{agg}}\) is the radius
of an equivalent volume sphere. Comparing to (Tomasko et al., 2008) all
the parameter ranges found in this work are narrower. Even without
combining the constraints from Blue and Red channels but selecting a
range with a minimum and maximum of both channels gives a narrower
monomer number range (see Figure 3) than 1500 - 6000 which was the
estimate by Tomasko et al. (2008). In this case, the smallest aggregate
has 2200 (from Blue Channel), the largest aggregate has 4700 (Red
channel) monomers (see table 2) with monomer size range between 0.038
(Red and Blue channel minimum) and 0.048 (Red channel maximum).

The refractive index ranges are shown in table 2.

\selectlanguage{english}
\FloatBarrier
\end{document}

